\documentclass[UKenglish,10pt,a4paper]{report}
\setlength\textwidth{145mm}
\setlength\textheight{247mm}
\setlength\oddsidemargin{5mm}
\setlength\evensidemargin{15mm}
\setlength\topmargin{0mm}
\setlength\headsep{0mm}
\setlength\headheight{0mm}
% \openright zařídí, aby následující text začínal na pravé straně knihy
\let\openright=\clearpage
\renewcommand{\baselinestretch}{1.0}%řádkování
\usepackage{booktabs}
\usepackage{multirow}
\usepackage{tabu}
\usepackage[hyphens]{url}
\usepackage{color}
\usepackage{listings}
\usepackage[utf8]{inputenc}
\usepackage{csquotes}
\usepackage[backend=biber]{biblatex}
\usepackage[pdftex]{graphicx}
\usepackage{epstopdf}
\usepackage{wrapfig}
\usepackage[section]{placeins}
\usepackage{amsfonts}
\usepackage{subcaption}
\usepackage{chronosys}
\usepackage{textcomp}
\usepackage{framed}
\usepackage{float}
\usepackage[czech, UKenglish]{babel}
\usepackage{multicol}
\usepackage{breakurl}
\usepackage[left]{eurosym}
\usepackage[binary-units]{siunitx}
\usepackage{makecell}
\usepackage{isodate}




\newcommand\ytl[2]{
\parbox[b]{8em}{\hfill{\color{cyan}\bfseries\sffamily #1}~$\cdots\cdots$~}\makebox[0pt][c]{$\bullet$}\vrule\quad \parbox[c]{7.5cm}{\vspace{7pt}\color{red!40!black!80}\raggedright\sffamily #2\\[7pt]}\\[-3pt]}

\usepackage[T1]{fontenc}
\usepackage{amsmath}
%\usepackage[toc,page]{appendix}
\usepackage[acronym,toc]{glossaries} 

\usepackage[unicode,breaklinks]{hyperref}  % Musí být za všemi ostatními balíčky
\usepackage{bookmark}




%\usepackage{pdfpages}

\addbibresource{ref.bib}


%% Pokud tiskneme oboustranně:
% \documentclass[12pt,a4paper,twoside,openright]{report}
% \setlength\textwidth{145mm}
% \setlength\textheight{247mm}
% \setlength\oddsidemargin{15mm}
% \setlength\evensidemargin{0mm}
% \setlength\topmargin{0mm}
% \setlength\headsep{0mm}
% \setlength\headheight{0mm}
% \let\openright=\cleardoublepage

%% Pokud používáte csLaTeX (doporučeno):
%\usepackage[czech]{babel}
%% Pokud nikoliv:

\begingroup
    \makeatletter
    \catcode`\-=\active
    \AtBeginDocument{
    \def\@@@cmidrule[#1-#2]#3#4{\global\@cmidla#1\relax
        \global\advance\@cmidla\m@ne
        \ifnum\@cmidla>0\global\let\@gtempa\@cmidrulea\else
        \global\let\@gtempa\@cmidruleb\fi
        \global\@cmidlb#2\relax
        \global\advance\@cmidlb-\@cmidla
        \global\@thisrulewidth=#3
        \@setrulekerning{#4}
        \ifnum\@lastruleclass=\z@\vskip \aboverulesep\fi
        \ifnum0=`{\fi}\@gtempa
        \noalign{\ifnum0=`}\fi\futurenonspacelet\@tempa\@xcmidrule}
    }
\endgroup


%% Použité kódování znaků: obvykle latin2, cp1250 nebo utf8:
%\usepackage[cp1250]{inputenc}

%\usepackage[T1]{fontenc} %tohle by mělo udělat Times New Roman 
%% Ostatní balíčky

%\usepackage{mathtools}
%\usepackage{amsthm}

%% Balíček hyperref, kterým jdou vyrábět klikací odkazy v PDF,
%% ale hlavně ho používáme k uložení metadat do PDF (včetně obsahu).
%% POZOR, nezapomeňte vyplnit jméno práce a autora.
\hypersetup{pdftitle=Optical communications}
\hypersetup{pdfauthor=Bc. et Bc. Jaroslav Svoboda, pdfkeywords={optical,network}}
%\hypersetup{hyperfootnotes=true}
\makeglossaries
%%% Drobné úpravy stylu

% Tato makra přesvědčují mírně ošklivým trikem LaTeX, aby hlavičky kapitol
% sázel příčetněji a nevynechával nad nimi spoustu místa. Směle ignorujte.
\makeatletter
\def\@makechapterhead#1{
  {\parindent \z@ \raggedright \normalfont
   \Huge\bfseries \thechapter. #1
   \par\nobreak
   \vskip 20\p@
}}
\def\@makeschapterhead#1{
  {\parindent \z@ \raggedright \normalfont
   \Huge\bfseries #1
   \par\nobreak
   \vskip 20\p@
}}
\makeatother

% Toto makro definuje kapitolu, která není očíslovaná, ale je uvedena v obsahu.
\def\chapwithtoc#1{
\chapter*{#1}}
\renewcommand{\familydefault}{\rmdefault}
\begin{document}
\selectlanguage{UKenglish}
% Trochu volnější nastavení dělení slov, než je default.
\lefthyphenmin=2
\righthyphenmin=2
\sisetup{range-phrase=--}
%%% Titulní strana práce
\begin{titlepage}


\pagestyle{empty}
\begin{center}

\large

TECHNICAL UNIVERSITY OF MADRID

\medskip

SCHOOL OF TELECOMMUNICATIONS SYSTEMS AND ENGINEERING

\vfill

{\bf\Large Semester Project}

\vfill
\centerline{\mbox{\includegraphics[width=60mm]{img/logo2.eps}\qquad{\includegraphics[width=60mm]{img/logo.eps}}}}
			
\vfill
\vspace{5mm}

{\LARGE Bc. et Bc. Jaroslav Svoboda}

{\LARGE Jaime Sancho}

{\LARGE Diallo Elhadj Sadou}

\vspace{10mm}

% Název práce přesně podle zadání
{\LARGE\bfseries } Smart Lighting

\vfill

% Název katedry nebo ústavu, kde byla práce oficiálně zadána
% (dle Organizační struktury MFF UK)
UBIQUITOUS AND SECURE NETWORKS AND SERVICES

\vfill

\begin{tabular}{rl}

Lecturer: & Jesus Rodriguez Molina  \\
\noalign{\vspace{2mm}}
Subject coordinator: & Jose Fernan Martinez Ortega \\
\noalign{\vspace{2mm}}
Degree programme: & Master in Systems and Services Engineering for the Information society \\
\end{tabular}

\vfill

% Zde doplňte rok
Madrid 2018

\end{center}
\end{titlepage}
\newpage

%%% Následuje vevázaný list -- kopie podepsaného "Zadání bakalářské práce".
%%% Toto zadání NENÍ součástí elektronické verze práce, nescanovat.


\openright

\noindent



%%% Strana s automaticky generovaným obsahem bakalářské práce. U matematických
%%% prací je přípustné, aby seznam tabulek a zkratek, existují-li, byl umístěn
%%% na začátku práce, místo na jejím konci.



\openright
\pagestyle{plain}
\setcounter{page}{1}
\tableofcontents



\chapter{Introduction and Executive Summary}
\chapter{Scheduling of workload}
\chapter{Requirement Analysis}
\begin{figure}[ht]
	\vspace{-2em}
	\caption{General use}
	\vspace{-1em}
	\centering
	\includegraphics[width=0.48\textwidth]{{"./img/Drawio/Use Cases/General use case diagram"}.png}
	\vspace{-2em}
\end{figure}
\begin{table}[h]
	\vspace{-1em}
	\centering
	\caption{General use}
	\label{tab:general_use}
	\begin{tabu}{|l|X[l]|}
		\hline
		\textbf{Name}        & \textbf{General use} \\ \hline
		Scope and Objectives & The lights on the lampposts switch on/off depending on the information of the movement and light sensor of the motes (if the light is low and the movement sensor detects something, the lights turn on, otherwise they turn off). The data collection centre receives periodically this data and it is capable to change some parameters or request some information of the motes. \\ \hline
		Actors               & \begin{itemize}
								\item Pedestrians/vehicles
								\item Data control department of the city
								\end{itemize}  \\ \hline
		Preconditions        & \begin{itemize}
								\item The sensor nodes are on
								\item A pedestrian/vehicle comes/goes away from the node
								\end{itemize}   \\ \hline
		Post-conditions      & \begin{itemize}
								\item The lights will turn on/off
								\item All the measured information will be sent to the city data manager
								\end{itemize} \\ \hline
		Sequence Description & \begin{enumerate}
								\item Obtain measurements of movement, temperature, humidity and light.
								\item If there is a change in the movement sensor and the light measured is low, toggle the lights on the node.
								\item Send the status of light to the data control department of the city.
								\end{enumerate} \\ \hline
		Exceptions           & The node is not able to communicate with the data control department of the city. \\ \hline
	\end{tabu}
		\vspace{-6em}
\end{table}



\renewcommand{\acronymname}{List of acronyms}

\printglossaries



%\bibliographystyle{plain} 
\printbibliography
%%% Přílohy k bakalářské práci, existují-li (různé dodatky jako výpisy programů,
%%% diagramy apod.). Každá příloha musí být alespoň jednou odkazována z vlastního
%%% textu práce. Přílohy se číslují.
\openright
\end{document}